% !TeX spellcheck = es_ES
%\documentclass[letterpaper,12pt]{article}
\documentclass[letterpaper,11pt]{article}

% Language & Input
  \usepackage[english]{babel} % Language: English
  \usepackage[ansinew]{inputenc} % Input
% Font
  \usepackage[T1]{fontenc} % Font encoding
  \usepackage{lmodern,microtype} % Typeface
% Math
  \usepackage{amsmath,amsthm,amssymb,amsfonts} % AMS math
  \usepackage{dsfont,mathrsfs,ushort} % Math style
% Style
  \usepackage[bf,sf,tiny,center]{titlesec} % Section titles 
  \usepackage{geometry} % Page & margins
  \usepackage{setspace} % Spacing
  \usepackage{enumitem} % Lists
  \usepackage{abstract} % Abstract
  \usepackage{color}
% References
  \usepackage[longnamesfirst]{natbib} % Manages bibiography
% Figures
  \usepackage{graphicx} % Additional options for includefigure
  \usepackage[font={small}]{caption}
% Tables
  \usepackage{booktabs} % Table format
  \usepackage{dcolumn} % Align on the decimal point of numbers in tabular columns.
  \usepackage{multirow} % Multirows in a table
  \usepackage{lscape} % Tables in landscape form
  \usepackage{ctable} % Tables with captions and notes all the same width 
% Hyperlinks
  \usepackage{hyperref}
% Addons
  \usepackage[l2tabu, orthodox]{nag} % Warn the user about the use of such obso­lete things

%%%%%%%%%%%%%%%%%%%%%%%%%%%%%%% Style & Structure
%% Titles (Carlos Perez)
%  \titlelabel{\thetitle.}
%  \titleformat*{\subsection}{\centering\sffamily\itshape}
%  \titleformat{\subsubsection}[runin]{\sffamily}{}{0em}{}[.]
%  \titlespacing{\subsubsection}{0pt}{*2}{\wordsep}
  
% Title page (Luis Cabral)
\makeatletter
    \renewcommand{\@maketitle}{%
    \newpage\parindent0in\null\vskip2em%
    {\LARGE\textbf{\textsf{\@title}}\par\vskip2em}%
    \@author\par\vskip2em%
    \@date%
    \par
    \thispagestyle{empty}\setcounter{page}{1}%\newpage
    }%
\makeatother
\renewenvironment{abstract}% Change the way abstract is presented
{\vskip2em\noindent{\textbf{\textsf{Abstract.}}}\hspace*{1ex}}{}

% Sections
  \renewcommand\thesection{\arabic{section}}
  \renewcommand\thesubsection{\Alph{subsection}}
% Bibliography
  \bibliographystyle{apa} %Alternative: {plain} {aea} {econometrica}
% Spacing & lists and margins
  \geometry{margin=2.54cm}
  \setstretch{1.15} % Set spacing
% Color
  \definecolor{ntr}{rgb}{0.8,0,0} % not too red
% Abstract
  \renewcommand{\abstractnamefont}{\sffamily\small\bfseries}
  \setlength{\abstitleskip}{-1em}
% Hyperlinks
  \hypersetup{colorlinks=true,linkcolor=ntr,citecolor=ntr,urlcolor=black,bookmarksnumbered=true}
% Notes to tables and figures
% Significant stars in tables
  %\newcommand{\sym}[1]{\rlap{#1}}
  \def\sym#1{\ifmmode^{#1}\else\(^{#1}\)\fi}
% Games
% Miscellaneous

%%%%%%%%%%%%%%%%%%%%%%%%%%%%%%% Theorems
  \theoremstyle{plain}
  \newtheorem{prop}{Proposition}
  \newtheorem{assumption} {Assumption}
  \newtheorem{theorem}{Theorem}
  \newtheorem{lem}{Lemma}
  \newtheorem{corollary}{Corollary}
  \theoremstyle{definition}
  \newtheorem{defi}{Definition}
  \newtheorem{ex}{Example}
  \theoremstyle{remark}
  \newtheorem{remark}{Remark}
  \def\propautorefname{Proposition}
  \def\propositionautorefname{Proposition}
  \def\defiautorefname{Definition}
  \def\lemautorefname{Lemma}
  \def\exampleautorefname{Example}
  \def\corollaryautorefname{Corollary}
  \def\exautorefname{Example}

\newcommand{\SR}{\mathbb{R}}

\DeclareMathOperator*{\Max}{m{a}x}
\DeclareMathOperator*{\Min}{m{i}n}
\DeclareMathOperator{\rank}{rank}
\DeclareMathOperator{\plim}{plim}
\DeclareMathOperator{\Var}{Var}
\DeclareMathOperator{\Avar}{AVar}
\DeclareMathOperator{\Cov}{Cov}
\DeclareMathOperator{\E}{\mathbb{E}}
\newcommand{\pto}{\stackrel{p}{\longrightarrow}}
\newcommand{\dto}{\stackrel{d}{\longrightarrow}}
\newcommand{\adistr}{\stackrel{a}{\sim}}
\newcommand{\epsi}{\varepsilon}
\newcommand{\tmle}{\hat \theta_{ML}}
\newcommand{\tgmm}{\hat \theta_{GMM}}
\newcommand\independent{\protect\mathpalette{\protect\independenT}{\perp}}
\def\independenT#1#2{\mathrel{\rlap{$#1#2$}\mkern2mu{#1#2}}}
\DeclareMathOperator{\N}{Normal}

\DeclareMathOperator{\Corr}{Corr}

\begin{document}

	\title {An\'alisis Econom\'etrico  \\ Tarea 4}
	\author{Mag\'ister en Econom\'ia, Universidad Alberto Hurtado}
	\date{Fecha de entrega: 5 de noviembre de 2026 (al principio de la ayudant\'ia)}

\maketitle

Total: 100 puntos

Los ejercicios 1, 2, 3 y 5 requieren estimar modelos con datos. Utilice R (u otro software
estad\'istico) y adjunte su c\'odigo. Reporte errores est\'andar robustos a
heteroscedasticidad y utilice un nivel de significancia del 5\% salvo que se indique otra
cosa. Los datos est\'an en el repositorio del curso, en formato R:
\href{https://github.com/ramirode/econometria_mae/raw/main/problem-sets/datos/veteran.rds}{\texttt{veteran.rds}},
\href{https://github.com/ramirode/econometria_mae/raw/main/problem-sets/datos/christensen_greene.rds}{\texttt{christensen\_greene.rds}},
\href{https://github.com/ramirode/econometria_mae/raw/main/problem-sets/datos/hprice2.rds}{\texttt{hprice2.rds}} y
\href{https://github.com/ramirode/econometria_mae/raw/main/problem-sets/datos/discrim.rds}{\texttt{discrim.rds}}. Se leen con \texttt{readRDS}.

\begin{enumerate}

\item (20 puntos) Utilice los datos de \texttt{veteran} del trabajo de Angrist (Econometrica,
    1998) sobre el efecto de haber servido en las fuerzas armadas de EE.UU. en los ingresos.
    Por razones de confidencialidad los datos est\'an agregados: cada fila es una celda de
    solicitantes con las mismas caracter\'isticas y \texttt{earnvar} es el ingreso anual promedio
    de la celda en 1988--1991 (en d\'olares). Trate cada celda como una observaci\'on, sin
    ponderadores: los promedios por grupo son promedios entre celdas. La
    variable \texttt{dvet} es igual a 1 si sirvi\'o en las fuerzas armadas (veterano) y
    \texttt{dnwhite} es igual a 1 si la persona no es de raza blanca.
    \begin{enumerate}
        \item (5 puntos) Proponga y estime un modelo de regresi\'on que permita diferencias de
            ingresos ($earnvar$) entre veteranos de raza blanca, veteranos de raza no blanca, no
            veteranos de raza blanca y no veteranos de raza no blanca. No utilice controles
            adicionales.
        \item (5 puntos) Para \emph{personas de raza blanca}, ?`cu\'al es la diferencia entre
            el ingreso medio de veteranos y de no veteranos? Para \emph{personas que no son de
            raza blanca}, ?`cu\'al es esa diferencia? Escriba ambas diferencias en funci\'on de
            los coeficientes del modelo.
        \item (5 puntos) Contraste si, para \emph{personas de raza blanca}, existen diferencias
            de ingresos entre veteranos y no veteranos. Realice un contraste similar para
            \emph{personas que no son de raza blanca}.
        \item (5 puntos) Contraste si la diferencia de ingresos entre veteranos y no veteranos
            es la misma para blancos y no blancos. ?`Qu\'e ventaja tiene estimar un solo modelo
            con dummies en lugar de dos modelos separados (uno para blancos y otro para no
            blancos)?
    \end{enumerate}

\item (20 puntos) Utilice los datos de \texttt{hprice2} del trabajo de Harrison y Rubinfeld
    (JEEM, 1978) sobre el efecto de la contaminaci\'on en el precio de las viviendas: 506
    comunidades del \'area metropolitana de Boston en 1970. Las variables son \texttt{price}
    (precio mediano de las viviendas, en d\'olares), \texttt{nox} (concentraci\'on de \'oxido de
    nitr\'ogeno en el aire, en partes por 100 millones), \texttt{rooms} (n\'umero promedio de
    habitaciones), \texttt{dist} (distancia ponderada a cinco centros de empleo) y
    \texttt{stratio} (ratio estudiantes por profesor). Considere el modelo
    \begin{align*}
    \log(price_i)=\beta_0+\beta_1 \log(nox_i)+\beta_2\, rooms_i +\beta_3 \log(dist_i) +\beta_4\, stratio_i +u_i.
    \end{align*}
    \begin{enumerate}
        \item (5 puntos) Contraste la hip\'otesis de que la elasticidad del precio de las
            viviendas con respecto a la contaminaci\'on es igual a $-1$.
        \item (5 puntos) Agregue $rooms^2$ al modelo y contraste la significatividad de
            $rooms^2$. Interprete el resultado del contraste.
        \item (5 puntos) En el modelo con $rooms^2$, escriba el efecto marginal de $rooms$
            sobre $\log(price)$, $\partial\log(price)/\partial rooms$, en funci\'on de $rooms$. ?`Para qu\'e
            n\'umero de habitaciones el efecto marginal pasa de negativo a positivo? ?`Cu\'antas
            observaciones de la muestra tienen menos habitaciones que ese valor? ?`Es
            preocupante que el efecto marginal estimado sea negativo en ese rango?
        \item (5 puntos) Construya un intervalo de confianza al 95\% para el efecto marginal
            de (c), evaluado en la media muestral de $rooms$. Compare con
            el intervalo de confianza para $\beta_2$ en el modelo sin $rooms^2$. ?`Qu\'e
            concluye?
    \end{enumerate}

\item (20 puntos) Los datos de \texttt{christensen\_greene} contienen costos, precios de
    insumos y producci\'on de empresas generadoras de electricidad de EE.UU. en 1955 y 1970. Los
    datos de 1955 son de Nerlove (1963), en la versi\'on que usan Christensen y Greene (JPE,
    1976). Las variables que utilizar\'a son:
    \begin{itemize}
        \item $year$: a\~no,
        \item $tc$: costo total,
        \item $q$: producci\'on,
        \item $pl$, $pk$ y $pf$: precios de la mano de obra, del capital y del combustible.
    \end{itemize}
    Suponga una funci\'on de producci\'on Cobb-Douglas
    \begin{align*}
    q_i = \exp(A+\eta_i)\, l_i^{\alpha_l} k_i^{\alpha_k} f_i^{\alpha_f},
    \end{align*}
    donde $\exp(A+\eta_i)$ es la productividad de la empresa, que se separa en la productividad
    promedio de la industria ($A$) y la productividad espec\'ifica de la empresa ($\eta_i$),
    $l_i$ es mano de obra, $k_i$ es capital f\'isico y $f_i$ es combustible.

    Si las empresas son tomadoras de precios en el mercado de los insumos y minimizan costos,
    se puede deducir la siguiente funci\'on de costos:
    \begin{align*}
    \log(tc_i) = \log(C_0) &+ \frac{\alpha_l}{\alpha_l+\alpha_k+\alpha_f} \, \log( pl_i) + \frac{\alpha_k}{\alpha_l+\alpha_k+\alpha_f} \, \log(pk_i) \\
    &+ \frac{\alpha_f}{\alpha_l+\alpha_k+\alpha_f} \, \log(pf_i) + \frac{1}{\alpha_l+\alpha_k+\alpha_f} \, \log(q_i) - \frac{1}{\alpha_l+\alpha_k+\alpha_f} \, \eta_i,
    \end{align*}
    donde $C_0$ es una constante que depende de $\alpha_l$, $\alpha_k$, $\alpha_f$ y $A$.

    \emph{Realice los siguientes ejercicios con los datos del a\~no 1955.}
    \begin{enumerate}
        \item (5 puntos) Estime la funci\'on de costos
        \begin{align*}
        \log(tc_i)=\beta_0+\beta_1 \log(pl_i)+\beta_2 \log(pk_i)+\beta_3 \log(pf_i)+\beta_4 \log(q_i)+u_i.
        \end{align*}
        Escriba $\beta_1$, $\beta_2$, $\beta_3$ y $\beta_4$ en funci\'on de $\alpha_l$,
        $\alpha_k$ y $\alpha_f$.
        \item (5 puntos) Muestre que la funci\'on de costos implica
            $\beta_1+\beta_2+\beta_3=1$ (homogeneidad de grado uno en precios: si todos los
            precios de los insumos se multiplican por $\lambda>0$, el costo total se
            multiplica por $\lambda$). Contraste esta restricci\'on: es una forma de contrastar
            con datos el modelo de minimizaci\'on de costos.
        \item (10 puntos) Utilice el coeficiente de $\log(q_i)$ para construir un intervalo de
            confianza al 95\% para los rendimientos a escala, $\alpha_l+\alpha_k+\alpha_f$.
            Contraste la hip\'otesis de rendimientos constantes a escala contra la alternativa
            bilateral.
    \end{enumerate}

\item (20 puntos) Suponga el modelo $y_i=\beta_0+\beta_1\,x_{1i}+\beta_2\,x_{2i}+u_i$. Bajo
    los supuestos usuales, el estimador MCO satisface $\sqrt{N}(\hat \beta-\beta)\dto\N(0,V)$ y
    $\hat V \pto V$. Usted quiere contrastar $H_0: \beta_1 = 0$ y $\beta_2 = 0$ contra la
    alternativa $H_1: \beta_1 \neq 0$ o $\beta_2 \neq 0$, con el estad\'istico de Wald
    \begin{align*}
    W_N=N\,(R'\hat \beta)'(R'\hat V R)^{-1}(R'\hat \beta), \qquad \text{donde} \qquad
    R'=\left(\begin{array}{ccc} 0 & 1 & 0\\ 0 & 0 & 1\end{array}\right).
    \end{align*}
    Demuestre que el estad\'istico de Wald se puede escribir como la siguiente funci\'on de los
    estad\'isticos t individuales ($t_{1,N}$, $t_{2,N}$) y del coeficiente de correlaci\'on
    estimado entre $\hat \beta_1$ y $\hat \beta_2$ ($\hat\rho$):
    \begin{align*}
    W_N = \frac{1}{1-\hat\rho^2}\left(t_{1,N}^2 + t_{2,N}^2-2 \hat \rho\, t_{1,N}\,t_{2,N}\right).
    \end{align*}
    Pista: Con $\widehat{\Avar}(\hat\beta)=\hat V/N$ y
    $se(\hat\beta_k)=[\widehat{\Avar}(\hat\beta_k)]^{1/2}$,
    \begin{align*}
    t_{1,N} = \frac{\hat \beta_1}{se(\hat \beta_1)}, \quad
    t_{2,N} = \frac{\hat \beta_2}{se(\hat \beta_2)} \quad \text{y} \quad
    \hat \rho =\frac{ \widehat{\Avar}(\hat \beta_1,\hat \beta_2)}{ se(\hat \beta_1)\, se(\hat \beta_2)}.
    \end{align*}

\item (20 puntos) Utilice los datos de \texttt{discrim} del trabajo de Graddy (JBES, 1997)
    sobre discriminaci\'on por raza e ingreso en los precios de las cadenas de comida r\'apida en
    New Jersey y Pennsylvania. Cada observaci\'on es un local de comida r\'apida, y las
    caracter\'isticas de la zona se miden a nivel del c\'odigo postal. Las variables son:
    \begin{itemize}
        \item $psoda$: precio de una bebida mediana en el local (en d\'olares),
        \item $prpblck$: proporci\'on de personas de raza negra en la zona,
        \item $income$: ingreso familiar mediano en la zona (en d\'olares),
        \item $prppov$: proporci\'on de personas pobres en la zona,
        \item $hseval$: valor mediano de la vivienda en la zona (en d\'olares).
    \end{itemize}
    Considere el modelo
    \begin{align*}
    \log(psoda_i)=\beta_0+\beta_1\, prpblck_i+\beta_2 \log(income_i) +\beta_3\, prppov_i +u_i.
    \end{align*}
    \begin{enumerate}
        \item (5 puntos) Estime el modelo. Interprete $\hat\beta_1$. ?`Es $prpblck$
            estad\'isticamente significativa al 5\%? ?`Y al 1\%?
        \item (5 puntos) Agregue $\log(hseval)$ al modelo. Interprete su coeficiente y
            compare $\hat\beta_1$ con el de (a).
        \item (5 puntos) En el modelo de (b), ?`son $\log(income)$ y $prppov$
            estad\'isticamente significativas en forma individual? Contraste si son
            conjuntamente significativas con el estad\'istico de Wald y reporte el p-valor.
        \item (5 puntos) Calcule la correlaci\'on muestral entre $\log(income)$ y $prppov$, y
            la correlaci\'on estimada entre $\hat\beta_2$ y $\hat\beta_3$, $\hat\rho$. ?`Qu\'e
            relaci\'on hay entre sus signos? Utilice la f\'ormula del ejercicio 4 (que vale para
            dos coeficientes cualesquiera de un modelo con m\'as regresores) para explicar el
            resultado de (c): compare $W_N$ con el valor que tendr\'ia si $\hat\rho=0$.
    \end{enumerate}

\end{enumerate}
\end{document}
